% ============================================================================
% PART I: CONSTRUCTION
% Section 1: Introduction
% ============================================================================
%
% Purpose
%   Frame the paper as a transition narrative: ordinary gauge QFT already
%   contains topological sectors and nonlocal observables; Chern--Simons theory
%   is where those structures stop being peripheral and become the defining
%   data; genuine TQFT is the formal completion; observables are the payoff.
%
% Target thesis (from tqft_observables_literature_review.md Sec. 11)
%   "Ordinary gauge QFT already contains topological sectors, characteristic
%    classes, and nonlocal observables, but Chern--Simons theory is the point
%    at which these structures cease to be peripheral and become the defining
%    data of the theory. The goal of this review is to explain that transition,
%    formalize it in the language of TQFT, and then show how the same structure
%    governs concrete observables ranging from Hall response and anyonic
%    braiding to the eta' mass, neutron EDM constraints, axion physics, and
%    monopole charge quantization."
%
% Status: TO BE WRITTEN (write last, after Sections 2--7)
%
% Length target: 3--5 pages
%
% Suggested subsections
%   1.1  What counts as "topological structure" inside ordinary QFT
%   1.2  Why Chern--Simons is the clean bridge to genuine TQFT
%   1.3  Observable families tracked by this paper
%   1.4  Roadmap (Part I construction, Part II physics organized by observable)
%
% Relevant sources to cite up front
%   Atiyah 1988 (\cite{Atiyah1988TQFT})                 -- functorial TQFT
%   Witten 1988 (\cite{Witten1988TQFT})                 -- cohomological TQFT
%   Witten 1989 (\cite{Witten1989Jones})                -- CS / Jones polynomial
%   Birmingham et al. 1991 (\cite{BirminghamBlauRakowskiThompson1991TopologicalFieldTheory})
%   Dunne 1999 (\cite{Dunne1999AspectsCS})              -- physics review of CS
%   Witten 2016 (\cite{Witten2016ThreeLecturesTopologicalPhases}) -- bridge
%   Wen 1995 (\cite{Wen1995TopologicalOrdersEdgeExcitations}) -- topological order
%   Stern 2008 (\cite{Stern2008AnyonsQHEReview})        -- observables review
%   Nayak et al. 2008 (\cite{NayakSimonSternFreedmanDasSarma2008NonabelianAnyons})
%
% Writing notes
%   - Do not open with an abstract TQFT definition.
%   - Open with the physics problem and Schwartz-level vocabulary.
%   - State the observable families up front so the reader knows the payoff.
%   - End with a visual map of Parts I and II.

\section{Introduction}
\label{sec:introduction}

% TODO after other sections are complete:
%   - Write thesis statement above as opening paragraph (Sec. 11 of literature review).
%   - Three-question framing: (i) what is topological structure inside ordinary QFT?
%     (ii) why does CS isolate it cleanly? (iii) how does this become formal TQFT?
%   - Observable families preview (transport, spectroscopic, braiding, defect).
%   - Part I / Part II roadmap.
